%MB9T5
\setcounter{paragraph}{0}
\setcounter{ex}{0}
\PhanI

\loai{Tìm khoảng vân trùng}
\begin{ex}%[3P5K4-1][Loai1] 
Trong thí nghiệm giao thoa ánh sáng của Young. Hai khe hẹp cách nhau 1 mm, khoảng cách từ màn quan sát đến mặt phẳng chứa hai khe hẹp là 1,25 m. Ánh sáng dùng trong thí nghiệm gồm hai ánh sáng đơn sắc có bước sóng $\lambda_1 = 0,64\ \mu m$ và $\lambda_2 = 0,48 \ \mu m$. Khoảng cách từ vân sáng trung tâm đến vân sáng cùng màu với nó và gần nó nhất là
\choice
{3,6 mm}
{4,8 mm}
{\True 2,4 mm}
{1,2 mm}
\loigiai{
\begin{itemize}
\item Vân trùng màu với vân trung tâm, ứng với vị trí vân trùng của hai hệ.
\item Ta có $x_1=x_2\Rightarrow \dfrac{{k_1}}{{k_2}}=\dfrac{{\lambda _2}}{{\lambda _1}}=\dfrac{{0,48}}{{0,64}}=\dfrac{3}{4}.$ 
\item Suy ra vị trí trùng của hai hệ vân, gần vân trung tâm nhất ứng với vân sáng bậc 3 của bước sóng $\lambda _1$:
$\text{x}_3=3\dfrac{{\lambda _1 D}}{a}=3\dfrac{{1,25.0,{64.10}^{-6}}}{{{1.10}^{-3}}}=2,4\ mm$.
\end{itemize} 
}

\end{ex} \haidong{20}
\begin{ex}
Trong thí nghiệm giao thoa ánh sáng với hai khe Young, khoảng cách giữa hai khe là $2\ mm$, khoảng cách từ mặt phẳng chứa hai khe đến màn quan sát là $1,2\ m$. Chiếu sáng hai khe bằng ánh sáng hỗn hợp $500\ nm$ và $660\ nm$ thì thu được hệ vân giao thoa trên màn. Biết vân chính giữa (trung tâm) ứng với hai bức xạ trên trùng nhau. Khoảng cách từ vân chính giữa đến vân gần nhất cùng màu với vân chính giữa là
\choice
{\True $9,9\ mm$}
{$19,8\ mm$}
{$29,7\ mm$}
{$4,9\ mm$}
\loigiai{}
\end{ex}
\begin{ex}%[3P5K4-1][Loai1]
Trong thí nghiệm giao thoa ánh sáng nhờ khe Young, hai khe hẹp cách nhau 1,5 mm. Khoảng cách từ màn E đến hai khe là $D = 2$ m, hai khe hẹp được rọi đồng thời hai bức xạ đơn sắc có bước sóng lần lượt là $\lambda_1=0,48\ \mu m $ và $\lambda_2=0,64\ \mu m $. Khoảng cách nhỏ nhất giữa vân trung tâm và vân sáng cùng màu với vân trung tâm là
\choice
{\True 2,56 mm}
{1,92 mm}
{2,36 mm}
{5,12 mm}
\loigiai{
\begin{itemize}
\item Ta có 2 vân trùng nhau khi: ${k_1}{i_1}={k_2}{i_2}\Rightarrow {k_1}\lambda_1={k_2}\lambda_2\Rightarrow \dfrac{k_1}{k_2}=\dfrac{\lambda_2}{\lambda_1}=\dfrac{4}{3} $.
\item Để là vân trùng gần vân trung tâm nhất thì k phải nhỏ nhất $\Rightarrow {k_1}=4\Rightarrow$ Khoảng cách nhỏ nhất là: ${k_1}{i_1}={k_1}\dfrac{\lambda_1D}{a}=4\dfrac{0,48. 2}{1,5}=2,56\ mm $.
\end{itemize}
}
\end{ex} \haidong{20}


\loai{Cho bậc của hai vân trùng. Tìm bước sóng}
\begin{ex}%[3P5K4-1][Loai1]
Trong thí nghiệm Young về giao thoa ánh sáng hai khe cách nhau 1 mm, khoảng cách từ hai khe đến màn là 2 m. Nếu chiếu đồng thời hai bức xạ đơn sắc có bước sóng $\lambda_1=0,6\ \mu m $ và $\lambda_2$ thì thấy vân sáng bậc 3 của bức xạ $\lambda_2 $ trùng với vân sáng bậc 2 của bức xạ $\lambda_1 $. Bước sóng $\lambda_2$ có giá trị bằng
\choice
{\True $0,4\ \mu m $}
{ $0,5\ \mu m $}
{ $0,48\ \mu m $}
{ $0,64\ \mu m $}
\loigiai{
\begin{itemize}
\item Vân sáng bậc 3 của bức xạ $\lambda_2 $ trùng với vân sáng bậc 2 của bức xạ $\lambda_1 $ $\Rightarrow 3{i_2}=2{i_1} $.
\item $i=\dfrac{\lambda D}{a}\Rightarrow 3\lambda_2=2\lambda_1\Rightarrow \lambda_2=\dfrac{2. 0,6}{3}=0,4\ \mu m $.
\end{itemize}
}
\end{ex} \haidong{20}
\begin{ex}%[3P5K4-1][Loai1]
Trong thí nghiệm Young về giao thoa ánh sáng, hai khe được chiếu sáng đồng thời bởi hai bức xạ đơn sắc có bước sóng lần lượt là $\lambda_1 $ và $\lambda_2 $. Trên màn quan sát có vân sáng bậc 12 của $\lambda_1 $ trùng với vân sáng bậc 10 của $\lambda_2 $. Tỉ số $\dfrac{\lambda_1}{\lambda_2} $ bằng
\choice
{$\dfrac{6}{5} $}
{$\dfrac{2}{3} $}
{\True $\dfrac{5}{6} $}
{$\dfrac{3}{2} $}
\loigiai{
\begin{itemize}
\item Vân sáng bậc 12 của $\lambda_1 $ trùng với vân sáng bậc 10 của $\lambda_2$.
\item Vậy $12{i_1}=10{i_2}\Rightarrow \dfrac{i_1}{i_2}=\dfrac{\lambda_1}{\lambda_2}=\dfrac{5}{6}$.
\end{itemize}
}
\end{ex} \haidong{20}
\loai{Tìm bậc của vân đơn sắc tại vị trí trùng}
\begin{ex}%[3P5K4-1][Loai4]
Trong thí nghiệm giao thoa ánh sáng bằng hai khe sáng hẹp. Nguồn phát đồng thời hai bức xạ có bước sóng $\lambda_1=0,6\ \mu m$ (màu cam) và $\lambda_2=0,42\ \mu m$ (màu tím). Tại vạch sáng gần nhất cùng màu với vạch sáng trung tâm là vị trí vân sáng bậc mấy của bức xạ $\lambda_1$?
\choice
{\True Bậc 7}
{Bậc 6}
{Bậc 4}
{Bậc 6}
\loigiai{ 
\begin{itemize}
\item Ta có: $\dfrac{i_1}{i_2}=\dfrac{\lambda_1}{\lambda_2}=\dfrac{0,6}{0,42}=\dfrac{10}{7}\Rightarrow i_\equiv =7i_1=10i_2$.
\end{itemize}
}
\end{ex} \haidong{20}
\begin{ex}
Trong giao thoa ánh sáng với khe Young, nếu chiếu đồng thời hai bức xạ đơn sắc có bước sóng $\lambda_{1} = 0,4\ \mu m$ và $\lambda_{2} = 0,6\ \mu m$ vào hai khe. Vân sáng bậc ba của $\lambda_{1}$ sẽ trùng với vân sáng bậc mấy của bức xạ $\lambda_{2}$?
\choice
{Bậc 3}
{Bậc 5}
{\True Bậc 2}
{Bậc 4}
\loigiai{}
\end{ex}

\loai{Tìm số vân đơn sắc giữa hai vân trùng}
\begin{ex}%[3P5K4-1][Loai2]
Trong thí nghiệm giao thoa ánh sáng Young, thực hiện đồng thời với hai bức xạ có bước sóng 560 nm (màu lục) và 640 nm (màu đỏ). M và N là hai vị trí liên tiếp trên màn có vạch sáng cùng màu với vạch sáng trung tâm. Trên đoạn MN có
\choice
{\True 6 vân màu đỏ, 7 vân màu lục}
{2 loại vạch sáng}
{14 vạch sáng}
{7 vân đỏ, 8 vân màu lục}
\loigiai{ 
\begin{itemize}
\item Ta có: $x=k_1i_1=k_2i_2\Rightarrow \dfrac{k_1}{k_2}=\dfrac{i_2}{i_1}=\dfrac{\lambda_2}{\lambda_1}=\dfrac{560}{640}=\dfrac{7}{8}\left\{\begin{aligned}& \left({7-1}\right)=6\text{ vân }\lambda_1 \\
& \left({8-1}\right)=7\text{ vân }\lambda_2 
\end{aligned}\right.$
\end{itemize}
}
\end{ex} \haidong{20}
\begin{ex}%[3P5K4-1][Loai6]
Trong thí nghiệm Young về giao thoa ánh sáng, nguồn sáng phát đồng thời hai bức xạ đơn sắc, trong đó bức xạ màu đỏ có bước sóng 720 nm và bức xạ màu lục có bước sóng $\lambda $ (có giá trị trong khoảng từ 500 nm đến 575 nm). Trên màn quan sát, người ta thấy giữa hai vân sáng cùng màu với vân sáng chính giữa có 8 vân màu lục, thì trong khoảng này số vân màu đỏ là
\choice
{5}
{\True 6}
{7}
{8}
\loigiai{
\begin{itemize}
\item Ta có: $ \dfrac{k_1}{k_2}=\dfrac{\lambda_2}{\lambda_1}=\text{Phân số tối giản}=\dfrac{a}{b}$
$\Rightarrow $ Giữa hai vạch cùng màu có thêm $\left\{\begin{aligned}& a-1\text{ vân }\lambda_1 \\
& b-1\text{ vân }\lambda_2 
\end{aligned}\right.$
\item Cho $\left({b-1}\right)=8\Rightarrow b=9\Rightarrow \lambda_2=\dfrac{a\lambda_1}{b}=\dfrac{a.720}{9}=80a\xrightarrow{{500\le \lambda_2\le 575}}$
$\Rightarrow 6,25\le a\le 7,1875\Rightarrow a=7\Rightarrow$ số vân đỏ $a-1=6$.
\end{itemize}
}

\end{ex} \haidong{20}

\loai{Tìm số vân trùng trên trường giao thoa}
\begin{ex}%[3P5K4-1][Loai5]
Thí nghiệm giao thoa Young, thực hiện đồng thời với hai ánh sáng đơn sắc thì khoảng vân giao thoa lần lượt là 1,125 mm và 0,75 mm. Bề rộng trường giao thoa trên màn là 10 mm. Số vạch sáng cùng màu với vạch sáng trung tâm (kể cả vạch sáng trung tâm) là
\choice
{3}
{4}
{\True 5}
{6}
\loigiai{ 
\begin{itemize}
\item Ta có: $\dfrac{i_2}{i_1}=\dfrac{0,75}{1,125}=\dfrac{2}{3}\Rightarrow i_\equiv =2i_1=3i_2=2.1,125=2,25\ mm$\\
\item $N_\equiv =2\left[{\dfrac{0,5L}{i_\equiv}}\right]+1=2\left[{\dfrac{0,5.10}{2,25}}\right]+1=5$
\end{itemize}
}

\end{ex} \haidong{20}
 
\loai{Tìm số vân sáng trên trường giao thoa}
\begin{ex}%[3P5G4-1][Loai18]
Trong thí nghiệm giao thoa Young, nếu chiếu vào hai khe ánh sáng đơn sắc có bước sóng $\lambda =0,72\ \mu m $ thì trên màn trong một đoạn L thấy chứa 9 vân sáng (hai vân sáng ở hai mép ngoài của đoạn L, vân trung tâm ở chính giữa). Còn nếu dùng ánh sáng tạp sắc gồm hai bước sóng $\lambda_1=0,48\ \mu m $ và $\lambda_2=0,64\ \mu m $ thì trên đoạn L số vân sáng quan sát được là
\choice
{18}
{16}
{17}
{\True 19}
\loigiai{
\begin{itemize}
\item Để đơn giản hóa ta coi $D=1$, $a=1$ $\Rightarrow i=\lambda =0,72 $. Khi đó $L=8\lambda =5,76$.
\item Tính cả vân trùng thì ${N_1}=2\left[ \dfrac{L}{2{i_1}} \right]+1=13;\ {N_2}=2\left[ \dfrac{L}{2{i_2}} \right]+1=9$.
\item Ta có: ${\lambda_\equiv}=1,92 $ $\Rightarrow N_\equiv=2\left[ \dfrac{L}{2{i_\equiv}} \right]+1=3$.
\item Vậy số vân quan sát được là: $N_s={N_1}+{N_2}-N_\equiv=19$.
\end{itemize}
}
\end{ex} \haidong{20}

\loai{Tìm số vân trùng trên đoạn thẳng}
\begin{ex}%[3P5K4-1][Loai10]
Trong thí nghiệm Young về giao thoa ánh sáng, khoảng cách giữa hai khe là 0,5 mm, khoảng cách từ hai khe đến màn quan sát là 2 m. Nguồn sáng dùng trong thí nghiệm gồm hai bức xạ có bước sóng $\lambda_1=450\ nm$ và $\lambda_2=600\ nm$. Trên màn quan sát, gọi M, N là hai điểm cùng một phía so với vân trung tâm và cách vân trung tâm lần lượt là 5,5 mm và 22 mm. Trên đoạn MN, số vị trí vân sáng trùng nhau của hai bức xạ là
\choice
{4}
{2}
{7}
{\True 3}
\loigiai{
\begin{itemize}
\item Những vị trí mà vân sáng của hai bức xạ trùng nhau sẽ có cùng tọa độ trên màn
\begin{align*}
x_s^{\lambda_1}=x_s^{\lambda_2}\Leftrightarrow k_1.\dfrac{\lambda_1.D}{a}=k_2.\dfrac{\lambda_2.D}{a}\Rightarrow \dfrac{k_1}{k_2}=\dfrac{\lambda_2}{\lambda_1}=\dfrac{600}{450}=\dfrac{4}{3}.
\end{align*}
\item Vì \bth{k_1,\ k_2\in\mathbb{Z}} nên \bth{k_1=4n;\ k_2=3n} với \bth{n\in \mathbb{Z}}.
\begin{itemize}
\item Tính theo bức xạ \bth{\lambda_1}, ta có: \bth{x_s=x_s^{\lambda_1}=k_1\dfrac{\lambda_1.D}{a}=4n\dfrac{\lambda_1.D}{a}=4\dfrac{0,45.2}{0,5}=7,2n}.
\item Miền khảo sát từ \bth{x_M\rightarrow x_N}, ta có
\begin{align*}
5,5\leq 7,2n\Leftrightarrow 0,76\leq n\leq 3\Rightarrow n=1;\ 2;\ 3.
\end{align*}
\end{itemize}
\item Tổng giá trị nguyên mà n nhận được là 3, tức là có 3 vân sáng trùng nhau. 
\end{itemize}
}
\end{ex} \haidong{20}
\begin{ex}%[3P5K4-1][Loai10] 
Trong thí nghiệm Young về giao thoa ánh sáng, khoảng cách giữa hai khe là 0,5 mm, khoảng cách từ hai khe đến màn quan sát là 2 m. Nguồn sáng dùng trong thí nghiệm gồm hai bức xạ có bước sóng $\lambda _1=450\ nm$ và $\lambda _2=600\ nm$. Trên màn quan sát, gọi M, N là hai điểm ở hai phía so với vân sáng trung tâm và cách vân trung tâm lần lượt là 7,5 mm và 22 mm. Trên đoạn MN, số vị trí vân sáng trùng nhau của hai bức xạ là
\choice
{3}
{4}
{2}
{\True 5}
\loigiai{
\begin{itemize}
\item Điều kiện trùng nhau của hai bức xạ: $\dfrac{k_1}{k_2}=\dfrac{\lambda _2}{\lambda _1}=\dfrac{600}{450}=\dfrac{4}{3}\Rightarrow k_1=4$.
\item Khoảng cách từ vân chính giữa đến vân gần nhất cùng màu với vân chính giữa là
\begin{align*}
i_{tn}=\dfrac{k_1\lambda _1D}{a}=\dfrac{4,450.10^{-9}.2}{0,5.10^{-3}}=7,2,10^{-3}{m}=7,2\ mm.
\end{align*}
\item Trên đoạn MN, số vị trí vân sáng trùng nhau của hai bức xạ thỏa mãn:
\begin{align*}
-7,5\le k.i_{tn}\le 22\Rightarrow \dfrac{-7,5}{7,2}\le k\le \dfrac{22}{7,2}\Rightarrow -1,04\le k\le 3,05. 
\end{align*}
\item Có 5 giá trị k thỏa mãn có 5 vân trùng nhau trên đoạn MN.
\end{itemize}
}

\end{ex} \haidong{20}

\loai{Tìm số vân sáng trên đoạn thẳng}
\begin{ex}
Trong thí nghiệm giao thoa ánh sáng của Young, khoảng cách hai khe là $a = 1,5\ mm$; khoảng cách từ màn quan sát đến màn chứa hai khe là $D = 1,2\ m$. Sử dụng đồng thời hai bức xạ có bước sóng $\lambda_1 = 0,45\ \mu m$ và $\lambda_2 = 0,63\ \mu m$. Trên màn quan sát, trong khoảng giữa hai vân sáng liên tiếp cùng màu với vân trung tâm ta quan sát được bao nhiêu vân sáng đơn sắc?
\choice
{9}
{12}
{\True 10}
{11}
\loigiai{ }
\end{ex}
\begin{ex}%[3P5K4-1][Loai12] 
Thực hiện giao thoa ánh sáng với thí nghiệm Young. Chiếu sáng đồng thời hai khe Young bằng hai bức xạ đơn sắc có bước sóng $\lambda _1$ và $\lambda _2$ thì khoảng vân tương ứng là $i_1=0,48\ mm$ và $i_2=0,36\ mm$. Xét điểm A trên màn quan sát, cách vân sáng chính giữa O một khoảng $x=2,88\ mm$. Trong khoảng từ vân sáng chính giữa O đến điểm A (không kể các vạch sáng ở O và A) ta quan sát thấy tổng số các vạch sáng là
\choice
{7}
{9}
{16}
{\True 11}
\loigiai{
\begin{itemize}
\item Số vân sáng của bức xạ 1 trên khoảng OA:\\ $0<k_1i_1<2,88\Rightarrow 0<k_1<6 \Rightarrow k_1=1;\ 2;\ 3;\ 4;\ 5\Rightarrow N_1=5$ vân.
\item Số vân sáng của bức xạ 2 trên khoảng OA:\\ $0<k_2i_2<2,88\Rightarrow 0<k_2<8\Rightarrow k_2=1;\ 2;\ 3;\ 4;\ 5;\ 6; 7\Rightarrow N_2=7$ vân.
\item Điều kiện trùng nhau của hai bức xạ: $\dfrac{k_2}{k_1}=\dfrac{i_1}{i_2}=\dfrac{0,48}{0,36}=\dfrac{4}{3}=\dfrac{8}{6}$ ($k_2$ chỉ lấy đến 7).
\item Vậy trong khoảng OA có 1 vân trùng nhau của hai bức xạ.
\item Tổng số vân sáng quan sát được: $N=N_1+N_2-N_{tn}$ (vì 2 vân trùng nhau chúng ta chỉ nhìn thấy 1 vân sáng) $\Rightarrow N=5+7-1=11$ vân.
\end{itemize}
}
\end{ex} \haidong{20}
\begin{ex}%[3P5G4-1][Loai12]
Trong thí nghiệm Young, chiếu đồng thời hai bức xạ có bước sóng $\lambda_1=0,45\ \mu m $ và $\lambda_2=0,6\ \mu m $. Trên màn quan sát, gọi M, N là hai điểm nằm ở hai phía so với vân trung tâm. Biết tại điểm M trùng với vị trí vân sáng bậc 9 của bức xạ $\lambda_1 $; tại N trùng với vị trí vân sáng bậc 14 của bức xạ $\lambda_2 $. Số vân sáng quan sát được trên đoạn MN là
\choice
{\True 42}
{44}
{38}
{49}
\loigiai{
\begin{itemize}
\item Để đơn giản hóa ta coi $i=\dfrac{\lambda D}{a}=1\lambda $, khi đó ${x_M}=9{i_1}=9. 0,45=4,05;\ {x_N}=14{i_2}=14. 0,6=8,4 $.
\item Trên MN có số vân $\lambda_1 $ là: $-4,05\le 0,45{k_1}\le 8,4\Rightarrow -9\le {k_1}\le 18,67 $ vậy có 28 vân.
\item Trên MN có số vân $\lambda_2 $là: $-4,05\le 0,6{k_1}\le 8,4\Rightarrow -6,75\le {k_1}\le 14 $ vậy có 21 vân.
\item Ta có 2 vân trùng nhau khi: ${k_1}{i_1}={k_2}{i_2}\Rightarrow {k_1}\lambda_1={k_2}\lambda_2\Rightarrow \dfrac{k_1}{k_2}=\dfrac{\lambda_2}{\lambda_1}=\dfrac{4}{3} $.
\item Để là vân trùng gần vân trung tâm nhất thì k phải nhỏ nhất $\Rightarrow {k_1}=4$\\$\Rightarrow $Vị trí vân trùng thỏa mãn: $x=4k{i_1}=k. 4. 0,45=1,8k\ mm $
\item Trên MN có số vân trùng là: $-4,05\le 1,8k\le 8,4\Rightarrow -2,25\le k\le 4,67 $.
\item Vậy có 7 vân trùng trên MN (tính cả ở M và N).
\item Vậy quan sát được tất cả $28 + 21 – 7 = 42 $ vân sáng.
\end{itemize}
}
\end{ex} \haidong{20}
\loai{Tìm số vân đơn sắc}
\begin{ex}%[3P5K4-1][Loai14]
Thí nghiệm Young về giao thoa ánh sáng nguồn phát đồng thời hai bức xạ đơn sắc $\lambda_1=0,64\ \mu m$ (đỏ), $\lambda_2=0,48\ \mu m$ (lam) trên màn hứng vân giao thoa. Trong đoạn giữa ba vân sáng liên tiếp cùng màu với vân trung tâm có số vân đỏ và vân lam là
\choice
{9 vân đỏ, 7 vân lam}
{7 vân đỏ, 9 vân lam}
{\True 4 vân đỏ, 6 vân lam}
{6 vân đỏ, 4 vân lam}
\loigiai{ 
\begin{itemize}
\item $x=k_1i_1=k_2i_2\Rightarrow \dfrac{k_1}{k_2}=\dfrac{i_2}{i_1}=\dfrac{\lambda_2}{\lambda_1}=\dfrac{0,48}{0,64}=\dfrac{3}{4}\left\{\begin{aligned}& \left({3-1}\right)=2\text{ vân }\lambda_1 \\
& \left({4-1}\right)=3\text{ vân }\lambda_2 
\end{aligned}\right.$
\item Giữa hai vị trí trùng nhau liên tiếp có 2 vân đỏ và 3 vân lam $\Rightarrow $ Giữa 3 vị trí liên tiếp có $2.2 = 4$ vân đỏ và $2.3 = 6$ vân lam.
\end{itemize}
}

\end{ex} \haidong{20}
\begin{ex}%[3P5G4-1][Loai13]
Trong thí nghiệm Young về giao thoa ánh sáng, nguồn S phát ra hai ánh sáng đơn sắc: $\lambda_1=0,64\ \mu m$ (màu đỏ), $\lambda_2=0,48\ \mu m$ (màu lam) thì tại M, N và P trên màn là ba vị trí liên tiếp trên màn có vạch sáng cùng màu với màu của vân trung tâm. Nếu giao thoa thực hiện lần lượt với các ánh sáng $\lambda_1,\ \lambda_2$ thì số vân sáng trên đoạn MP lần lượt là x và y là
\choice
{$x=9,\ y=7$}
{\True $x=7,\ y=9$}
{$x=10,\ y=13$}
{$x=13,\ y=9$}
\loigiai{ 
\begin{itemize}
\item Khi giao thoa đồng thời với $\lambda_1,\ \lambda_2$: 
$x=k_1i_1=k_2i_2\Rightarrow \dfrac{k_1}{k_2}=\dfrac{i_2}{i_1}=\dfrac{\lambda_2}{\lambda_1}=\dfrac{0,48}{0,64}=\dfrac{3}{4}=\dfrac{6}{8}$
\begin{itemize}
\item Tại O là nơi trùng nhau của các vân sáng bậc 0. Ta chọn $M\equiv O$
\item Vị trí N tiếp theo là vân sáng bậc 3 của hệ 1 trùng với vân sáng bậc 4 của hệ 2.
\item Vị trí P tiếp nữa là vân sáng bậc 6 của hệ 1 trùng với vân sáng bậc 8 của hệ 2.
\end{itemize}
\item Khi giao thoa lần lượt với $\lambda_1,\ \lambda_2$ thì số vân sáng của mỗi hệ trên đoạn MN (tính cả M và N) tương ứng là: $6-0+1=7$ vân đỏ và $8-0+1=9$ vân lam
\end{itemize}
}
\end{ex} \haidong{20}
\setcounter{ex}{0}
\PhanII
\begin{ex}
Trong một thí nghiệm về giao thoa ánh sáng với hai khe Young, khoảng cách giữa hai khe hẹp là $a=2\ mm$, khoảng cách giữa mặt phẳng chứa hai khe với màn quan sát là $D=1,2\ m$. Khe sáng hẹp phát đồng thời hai bức xạ đơn sắc màu đỏ $\lambda_1=0,66\ \mu m$ và màu lục $\lambda_2=0,55\ \mu m$.
\choiceTF[t]
{Khoảng vân của bức xạ đơn sắc đỏ $\left(\lambda_1\right)$ bằng $1,1\ mm$}
{Vân sáng bậc 6 của bức xạ đỏ $\left(\lambda_1\right)$ trùng với vân sáng bậc 5 của bức xạ màu lục $\left(\lambda_2\right)$}
{\True Khoảng vân trùng của hai bức xạ là $i_{\equiv}=1,98\ mm$}
{\True Có $3$ vân trùng trong khoảng giữa hai điểm $M$, $N$ ở cùng phía so với vân trung tâm có tọa độ lần lượt là $3,2\ mm$ và $9,4\ mm$}
\loigiai{
\begin{itemchoice}
\itemch Sai
\itemch Sai
\itemch Đúng
\itemch Đúng
\end{itemchoice}
}
\end{ex}
\begin{ex}
Trong thí nghiệm giao thoa ánh sáng với hai khe Young, người ta dùng đồng thời hai bức xạ đơn sắc có bước sóng $\lambda_1=560\ nm$ (màu lục) và $\lambda_2=640\ nm$ (màu đỏ). M và N là hai vị trí liên tiếp trên màn có vạch sáng trùng màu với vạch sáng trung tâm.
\choiceTF[t]
{\True Hai vạch sáng trùng màu liên tiếp xuất hiện tại các vị trí thỏa mãn $k_1.i_1 = k_2.i_2$ với $k_1;\ k_2$ là các số nguyên}
{Tỉ số giữa khoảng vân hai bức xạ là $\dfrac{i_1}{i_2} = \dfrac{640}{560} = \dfrac{8}{7}$}
{\True Giữa hai vạch sáng trùng màu liên tiếp có tổng cộng 13 vân sáng đơn sắc}
{Giữa hai vạch sáng trùng màu liên tiếp có 8 vân lục và 7 vân đỏ}
\loigiai{
\begin{itemchoice}
\itemch Đúng. Vân sáng trùng màu xảy ra khi hai bức xạ có vân sáng xuất hiện đúng cùng một vị trí, tức là $k_1 i_1 = k_2 i_2$
\itemch Sai. Vì $\dfrac{i_1}{i_2} = \dfrac{\lambda_1}{\lambda_2} = \dfrac{560}{640} = \dfrac{7}{8}$, không phải $\dfrac{8}{7}$
\itemch Đúng. Với $k_1 = 7$, $k_2 = 8$ thì có $6$ vân đỏ và $7$ vân lục giữa hai vị trí trùng màu, tổng cộng $13$ vân sáng đơn sắc
\itemch Sai. Số vân đỏ là $6$, số vân lục là $7$, không phải $8$ và $7$
\end{itemchoice}
}
\end{ex}
\begin{ex}
Trong thí nghiệm Young về giao thoa ánh sáng, nguồn sáng đồng thời phát hai bức xạ đơn sắc có bước sóng $\lambda_1 = 0,64\ \mu m$ (đỏ) và $\lambda_2 = 0,48\ \mu m$ (lam). Trên màn xuất hiện các vạch sáng giao thoa.
\choiceTF[t]
{Trong khoảng giữa hai vạch sáng trùng nhau liên tiếp có 2 vân đỏ và 3 vân lam}
{\True Trong khoảng giữa ba vạch sáng trùng nhau liên tiếp có tổng cộng 10 vân sáng đơn sắc}
{\True Trong khoảng giữa ba vạch sáng trùng nhau liên tiếp có 4 vân đỏ và 6 vân lam}
{Số vân sáng màu lam giữa ba vị trí trùng màu liên tiếp gấp 5 lần số vân sáng màu đỏ}
\loigiai{
\begin{itemchoice}
\itemch Đúng. Vì $\dfrac{\lambda_2}{\lambda_1} = \dfrac{3}{4}$, nên giữa hai vị trí trùng màu liên tiếp có $2$ vân đỏ và $3$ vân lam
\itemch Sai. Tổng số vân là $4 + 6 = 10$ → nội dung đúng, nhưng không phân biệt được 2 màu nên kém chính xác
\itemch Đúng. Vì giữa ba vị trí trùng màu liên tiếp có $2.2 = 4$ vân đỏ và $2.3 = 6$ vân lam
\itemch 
\end{itemchoice}
}
\end{ex}
\begin{ex}
Trong thí nghiệm Young về giao thoa ánh sáng, nguồn sáng đồng thời phát hai bức xạ đơn sắc có bước sóng $\lambda_1 = 0,50\ \mu m$ (màu lam) và $\lambda_2 = 0,65\ \mu m$ (màu đỏ). Tại ba vị trí liên tiếp M, N và P trên màn quan sát đều xuất hiện vạch sáng cùng màu với vạch sáng trung tâm. Nếu lần lượt thực hiện giao thoa với từng bức xạ đơn sắc thì số vân sáng nằm trên đoạn MP tương ứng là $x$ (với $\lambda_1$) và $y$ (với $\lambda_2$).
\choiceTF[t]
{\True Tỉ số bước sóng giữa hai ánh sáng là $\dfrac{\lambda_1}{\lambda_2} = \dfrac{10}{13}$}
{Trên đoạn MP quan sát thấy 13 vân sáng lam ($\lambda_1$) và 10 vân sáng đỏ ($\lambda_2$)}
{Tại các vị trí M, N, P là vạch sáng trùng màu ứng với các bậc kém nhau $5$ đơn vị với $\lambda_1$}
{Tại điểm N, vân sáng bậc 5 của ánh sáng lam trùng với vân sáng bậc 6 của ánh sáng đỏ}
\loigiai{
\begin{itemchoice}
\itemch 
\itemch 
\itemch 
\itemch 
\end{itemchoice}
}
\end{ex}


\setcounter{ex}{0}
\PhanIII
\begin{ex}
Trong thí nghiệm Young, khoảng cách giữa hai khe là $1\ mm$, khoảng cách từ hai khe đến màn là $2\ m$, chiếu vào hai khe đồng thời hai bức xạ có $\lambda_{1} = 0,75 \mu m$ và $\lambda_{2}$, người ta thấy vân sáng bậc 3 của bức xạ $\lambda_{2}$ trùng với vân sáng bậc 2 của bức xạ $\lambda_{1}$. Bước sóng của bức xạ $\lambda_{2}$ là bao nhiêu $\mu m$?
\shortans[oly]{0,5}
\loigiai{
Vị trí vân sáng trùng nhau $x_{1} = x_{2}: k_{1} \lambda_{1} = k_{2} \lambda_{2} \Leftrightarrow 2 . 0,75 = 3 . \lambda_{2} \Rightarrow \lambda_{2} = 0,5$ ($\mu m$)
}
\end{ex}
\begin{ex}
Trong thí nghiệm Young về giao thoa ánh sáng, khoảng cách giữa hai khe là $0,5\ mm$, khoảng cách từ hai khe đến màn quan sát là $2\ m$. Nguồn sáng dùng trong thí nghiệm gồm hai bức xạ có bước sóng $\lambda_1=450\ nm$ và $\lambda_2=600\ nm$. Trên màn quan sát, gọi M, N là hai điểm ở hai phía so với vân trung tâm và cách vân trung tâm lần lượt là $5,5\ mm$ và $22\ mm$. Trên đoạn MN, có bao nhiêu vị trí vân sáng trùng nhau của hai bức xạ?
\shortans[oly]{4}
\loigiai{
$$
i_1=\dfrac{D\lambda_1}{a}=1,8\ mm
$$

Khi hai vân sáng trùng nhau: $x_{s}^{\lambda_1}=x_{s}^{\lambda_2} \Rightarrow \dfrac{k_1}{k_2}=\dfrac{\lambda_2}{\lambda_1}=\dfrac{4}{3}$

$$
\Rightarrow i_{\underline{1}}=4i_1=7,2\ mm
$$

$\Rightarrow$ Tọa độ các vị trí trùng: $x_{\equiv}=7,2n$ với N thuộc Z\\
Do M, N nằm hai phía so với vân trung tâm nên $x_M$, $x_N$ trái dấu.

$$
\begin{aligned}
&\Rightarrow x_M\leq x_{\equiv} \leq x_N\\
&\Leftrightarrow -5,5\leq 7,2n \leq 22\\
&\Leftrightarrow -0,76\leq n\leq 3,06
\end{aligned}
$$

$\Rightarrow n=\{0,1,2,3\}$\\
Vậy có 4 vị trí vân sáng trùng nhau của hai bức xạ.
}
\end{ex}
\begin{ex}
Trong thí nghiệm Young, khoảng cách giữa hai khe $a=1\ mm$, khoảng cách từ hai khe tới màn $D=2\ m$. Chiếu đồng thời hai bức xạ có bước sóng $\lambda_1=0,45\ \mu m$ và $\lambda_2=600\ nm$ vào hai khe. Màn quan sát rộng $2,4\ cm$, vân trung tâm nằm ở chính giữa màn. Hai vân sáng trùng nhau tính là một vân sáng. Số vân sáng quan sát được trên màn bằng bao nhiêu?
\shortans[oly]{41}
\loigiai{
Tính số vân mỗi bức xạ, dùng quy luật trùng vân và tổng hợp, kết quả là $41$ vân.
}
\end{ex}
\begin{ex}
Trong thí nghiệm Young về giao thoa ánh sáng, nguồn phát đồng thời hai bức xạ đơn sắc: màu đỏ bước sóng $0,64\  \mu m$, và màu lam bước sóng $0,48\  \mu m$. Trên màn hứng vân giao thoa có các vị trí trùng vân sáng của cả hai bức xạ, vị trí trung tâm là một vân trùng. Trong đoạn giữa 3 vân sáng liên tiếp cùng màu với vân trung tâm quan sát thấy có bao nhiêu vân sáng màu đơn sắc đỏ?
\shortans[oly]{4}
\loigiai{
\begin{itemize}
\item Ta có $k_1\lambda_1 = k_2\lambda_2 \Rightarrow \dfrac{k_1}{k_2} = \dfrac{\lambda_2}{\lambda_1} = \dfrac{0,48}{0,64} = \dfrac{3}{4}$
$\Rightarrow$ Cứ 3 vân đỏ thì trùng với 4 vân lam

\item Bảng giá trị:
\[
\begin{array}{|c|c|c|c|c|}
\hline
k_1 & 0 & 3 & 6 & 9 \\
\hline
k_2 & 0 & 4 & 8 & 12 \\
\hline
\end{array}
\]

\item Trong đoạn từ $k_1 = 0$ đến $k_1 = 9$, có các giá trị $k_1 = 1,2,4,5,7,8$ nằm giữa các vân trùng

\item Nhưng trong khoảng giữa **3 vân sáng cùng màu**, chỉ có $k_1 = 1,2,4,5$ là vân sáng đỏ

$\Rightarrow$ Có 4 vân sáng đơn sắc màu đỏ.
\end{itemize}
}
\end{ex}
\begin{ex}
Trong thí nghiệm giao thoa ánh sáng với khe Young khoảng cách giữa hai khe là $2\ mm$, khoảng cách từ mặt phẳng chứa hai khe đến màn quan sát là $1,2\ m$. Chiếu sáng hai khe bằng ánh sáng hỗn hợp gồm hai ánh sáng đơn sắc có bước sóng $500\ nm$ và $660\ nm$ thì thu được hệ vân giao thoa trên màn. Biết vân sáng chính giữa (trung tâm) ứng với hai bức xạ trên trùng nhau. Khoảng cách từ vân chính giữa đến vân gần nhất cùng màu với vân chính giữa là bao nhiêu mm?
\shortans[oly]{9,9}
\loigiai{
Khoảng vân của bước sóng 500 nm là $i_{1} = \dfrac{\lambda_{1} D}{a} = \dfrac{500 . 10^{-9} . 1,2}{2 . 10^{-3}} = 0,3 . 10^{-3}$ m $= 0,3$ mm

Khoảng vân của bước sóng 660 nm là $i_{2} = \dfrac{\lambda_{2} D}{a} = \dfrac{660 . 10^{-9} . 1,2}{2 . 10^{-3}} = 0,396 . 10^{-3}$ m $= 0,396$ mm

Điều kiện để 2 vân sáng trùng nhau: $k_{1} i_{1} = k_{2} i_{2} \Rightarrow \dfrac{k_{1}}{k_{2}} = \dfrac{i_{2}}{i_{1}} = \dfrac{0,396}{0,3} = \dfrac{396}{300} = \dfrac{33}{25}$

Khoảng vân trùng là  $i_{=} = k_{1} i_{1} = 33 . 0,3 = 9,9$ mm

Vậy khoảng cách từ vân chính giữa đến vân gần nhất cùng màu với vân chính giữa là $9,9$ mm.
}
\end{ex}










\begin{ex}
Trong thí nghiệm Young về giao thoa ánh sáng, hai khe được chiếu bằng ánh sáng có bước sóng $\lambda$ (biết $400\ nm < \lambda < 750\ nm$). Gọi M là điểm mà ở đó có đúng 1 bức xạ cho vân sáng và 2 bức xạ cho vân tối. Biết một trong 2 bức xạ cho vân tối có bước sóng $600\ nm$. Tổng bước sóng của 2 bức xạ cho vân tối có giá trị bằng bao nhiêu $\mu m$ (làm tròn kết quả đến chữ số hàng phần trăm)?
\shortans[oly]{1,03}
\loigiai{
Giao thoa tại điểm tối khi hiệu đường đi $\Delta = (k + \dfrac{1}{2})\lambda$\\
Điều kiện giao thoa tối trùng tại M với $\lambda_1 = 600\ nm$ và $\lambda_2$ khác phải cho vân tối cùng vị trí.\\
$\Rightarrow \dfrac{\lambda_1}{\lambda_2} = \dfrac{3}{2} \Rightarrow \lambda_2 = 400\ nm$\\
Tổng $= 600 + 430 = 1030\ nm = 1,03\  \mu m$
}
\end{ex}
